% =============================================================================
% 2.7 Workflow de traitement
% =============================================================================
\section{Workflow de traitement}
\label{sec:ch2-workflow}

\subsection{Processus d'ingestion}

Le processus d'ingestion, porté par l'horloge A, est détaillé à la
figure~\ref{fig:ingestion}.

\begin{figure}[htbp]
  \centering
  \resizebox{\linewidth}{!}{% Chaîne d'ingestion documentaire
\begin{tikzpicture}[node distance=5mm and 9mm]

  \node[btpboxstrong, text width=2.1cm] (up) {Dépôt du\\fichier};
  \node[btpdb, right=of up, minimum width=2.0cm] (obj) {Stockage\\objet};
  \node[btpbox, text width=2.3cm, right=of obj] (detect)
    {Détection\\du type\\[-1pt]\btplabel{numérique / numérisé}};

  \node[btpbox, text width=2.5cm, above right=4mm and 9mm of detect] (parse1)
    {Analyse\\structurelle\\[-1pt]\btplabel{PDF numérique}};
  \node[btpbox, text width=2.5cm, below right=4mm and 9mm of detect] (parse2)
    {Reconnaissance\\optique\\[-1pt]\btplabel{document numérisé}};

  \node[btpbox, text width=2.4cm, right=9mm of {$(parse1)!0.5!(parse2)$}] (chunk)
    {Découpage\\hiérarchique};
  \node[btpboxalt, text width=2.3cm, below=14mm of chunk] (embed)
    {Calcul des\\plongements};
  \node[btpdb, left=12mm of embed, minimum width=2.1cm] (vec) {Index\\vectoriel};
  \node[btpdb, right=12mm of embed, minimum width=2.3cm] (sql) {Base\\relationnelle};

  \draw[btparrow] (up) -- (obj);
  \draw[btparrow] (obj) -- (detect);
  \draw[btparrow] (detect) -- (parse1);
  \draw[btparrow] (detect) -- (parse2);
  \draw[btparrow] (parse1) -- (chunk);
  \draw[btparrow] (parse2) -- (chunk);
  \draw[btparrow] (chunk) -- (embed);
  \draw[btparrow] (embed) -- (vec);
  \draw[btparrow] (embed) -- (sql);

  \node[btplabel, below=4mm of embed, text width=9.5cm]
    {Le statut du document passe de « en attente » à « prêt » seulement
     lorsque l'indexation est confirmée dans les deux magasins};

\end{tikzpicture}
}
  \caption{Chaîne d'ingestion documentaire}
  \label{fig:ingestion}
\end{figure}

Cinq étapes s'enchaînent.

\begin{enumerate}
  \item \btpterm{Dépôt.} L'utilisateur transmet un fichier. Le système
        enregistre le binaire dans le stockage objet, crée un
        enregistrement « en attente » et place un travail dans la file.
        La réponse HTTP est immédiate.
  \item \btpterm{Analyse.} Un travailleur détecte si le document est
        numérique ou numérisé, puis applique l'analyse structurelle ou la
        reconnaissance optique avec préservation de la mise en page.
  \item \btpterm{Découpage.} Les titres et articles structurent le
        découpage parent-enfant. Un préfixe de métadonnées (type de pièce,
        page, identifiant) est injecté dans le texte indexé, afin que la
        recherche lexicale bénéficie de ces ancres.
  \item \btpterm{Indexation.} Les plongements sont calculés et écrits dans
        l'index vectoriel~; le texte et les métadonnées le sont dans la
        base relationnelle. Le statut passe à « prêt » seulement après
        confirmation des deux écritures.
  \item \btpterm{Graphe (phase ultérieure).} Un travail distinct extraire
        entités et relations, fusionnées dans le magasin graphe.
\end{enumerate}

\subsection{Processus de génération}

Le processus de génération, porté par l'horloge B, est représenté à la
figure~\ref{fig:workflow-generation}.

\begin{figure}[htbp]
  \centering
  \resizebox{0.92\linewidth}{!}{% Workflow de génération d'une section avec supervision humaine
\begin{tikzpicture}[node distance=5.5mm and 12mm]

  \node[btpboxstrong, text width=2.6cm] (s1) {1. Routage de\\l'intention};
  \node[btpbox, text width=2.6cm, below=of s1] (s2) {2. Construction\\de la requête};
  \node[btpbox, text width=2.6cm, below=of s2] (s3) {3. Recherche\\hybride};
  \node[btpbox, text width=2.6cm, below=of s3] (s4) {4. Assemblage\\du contexte};
  \node[btpboxalt, text width=2.6cm, below=of s4] (s5) {5. Génération\\structurée};

  \node[btpdecision, text width=2.3cm, right=16mm of s5] (v)
    {Validation\\automatique};

  \node[btpbox, text width=2.7cm, above=12mm of v] (retry)
    {Reprise bornée\\[-1pt]\btplabel{2 tentatives au plus}};

  \node[btpboxstrong, text width=2.7cm, above=12mm of retry, fill=btpsky] (hitl)
    {6. Revue humaine\\[-1pt]\btplabel{approuver / corriger / relancer}};

  \node[btpbox, text width=2.7cm, above=12mm of hitl] (save)
    {7. Enregistrement\\de la version};

  \foreach \i/\j in {1/2, 2/3, 3/4, 4/5} {\draw[btparrow] (s\i) -- (s\j);}
  \draw[btparrow] (s5) -- (v);
  \draw[btparrow] (v) -- node[btplabel, right, xshift=1mm] {échec} (retry);
  \draw[btparrow] (retry) -- node[btplabel, right, xshift=1mm] {succès} (hitl);
  \draw[btparrow] (v.east) -- ++(0.7,0) |- node[btplabel, pos=0.25, right] {succès} (hitl.east);
  \draw[btparrow] (hitl) -- (save);
  \draw[btparrow] (retry.west) -- ++(-0.6,0) |- (s5.west);

  \node[btplabel, below=3mm of s5, text width=4.4cm]
    {Sortie contrainte par un schéma~: blocs, citations, exigences couvertes,
     abstentions explicites};

\end{tikzpicture}
}
  \caption{Workflow de génération d'une section avec supervision humaine}
  \label{fig:workflow-generation}
\end{figure}

La validation automatique qui précède la revue humaine n'est pas un
second modèle juge. Ce sont des contrôles déterministes~: le document
produit est-il conforme au schéma~? Les identifiants cités existent-ils~?
Aucun titre de premier ou deuxième niveau n'a-t-il été inventé hors plan~?
Tout fait chiffré est-il accompagné d'une citation~? Un échec déclenche
au plus deux reprises~; au-delà, la section est présentée à l'humain avec
le diagnostic.

\subsection{Routage des intentions}

L'entrée conversationnelle de l'espace de travail n'est pas un dialogue
libre. Un routeur, combinant règles et modèle, classe l'intention avant
tout appel coûteux.

\begin{btptable}{Routage des intentions}
\label{tab:intentions}
\begin{tabularx}{\linewidth}{@{}L{2.6cm} Y L{3.4cm} L{2.4cm}@{}}
\toprule
\rowcolor{btppale}
\btpth{Intention} & \btpth{Indices} & \btpth{Traitement} & \btpth{Sortie} \\
\midrule
Export & « exporter », « télécharger », « pdf » & Assemblage immédiat des sections approuvées. & Fichier \\
Générer une section & « génère », « rédige », « section » & Recherche, génération, persistance, revue. & Section \\
Question documentaire & Autre question ancrée dans le dossier & Recherche hybride et réponse citée. & Réponse \\
Échange social & Salutation, absence d'ancrage & Dialogue court, sans recherche. & Réponse \\
\bottomrule
\end{tabularx}
\end{btptable}

Ce routage évite deux écueils~: lancer une génération coûteuse pour une
simple salutation, et lancer un export alors qu'aucune section n'est
encore approuvée. Dans le second cas, le système refuse plutôt que de
produire un livrable incomplet.
